\documentclass[11pt]{article}
\usepackage{fontspec}
\usepackage[french,english]{babel}
\usepackage{amsmath}
\usepackage{amssymb}
\usepackage{amsthm}
\usepackage{mathrsfs}
\usepackage{tikz-cd}
\usepackage{enumitem}
\usetikzlibrary{babel}
\newcommand{\printedpage}[1]{\marginpar{\footnotesize p.~#1}}
\newcommand{\msp}[1]{\marginpar{\footnotesize [#1]}}
\newcommand{\Hom}{\operatorname{Hom}}
\newcommand{\Ob}{\operatorname{Ob}}
\newcommand{\bD}{\mathbf{D}}
\newcommand{\bDelta}{\boldsymbol{\Delta}}
\begin{document}

\printedpage{34}
% head: A. Grothendieck : Dérivateurs, Ch. II, Édition des Universités de Montpellier II et Paris VII
This shows the equivalence of d, d$'$, which can also be seen from the equivalence of each of these two conditions with d$''$, by the cofinality criterion of the preceding paragraph. Let D be this condition on $u$. By (4.8) we have D $\Longrightarrow$ A, hence % [sic] l'equivalence -> l'équivalence, accent missing; précedent -> précédent, accent missing
\[
\mathrm{D} \Longrightarrow \mathrm{A} \Longrightarrow \mathrm{c},
\]
it remains to prove that c $\Longrightarrow$ D, more precisely c $\Longrightarrow$ d$''$ and c$'$ $\Longrightarrow$ d$''$ \texttt{[c$'$?]}. Let us make explicit the adjunction morphism (for $i$ in $I$)
\[
u_!\underbrace{u^*i} \longrightarrow i,
\]
i.e. for $j$ in $I$
\[
\underbrace{(u_!u^*i)(j)}_{\pi_0(j\backslash(I'/u^*(i)))} \longrightarrow i(j) = \Hom(j,i).
\]

\msp{page 54}Now $I'/u^*(i) \simeq I'/i$, so we have the arrow
\[
(u_!u^*(i))(j) \simeq \pi_0(j\backslash(I'/i)) \longrightarrow \Hom_I(j,i). \tag*{(4.10)}
\]
Now we see at once that for $\alpha : j \longrightarrow i$ in $\Hom(j,i)$, the fiber of (4.10) at $u$ identifies with $\pi_0(j\backslash\underset{\alpha}{I'}/i\ )$, more precisely, $j\backslash\underset{\alpha}{I'}/i$\ \ is the open and closed subcategory of $J = j\backslash(I'/i)$ corresponding to the part of $\pi_0(J)$ which lies over $\alpha$ in (4.10). This completes the proof of 4.1.

NB Another way to prove that c $\Longrightarrow$ A is to note that the two sides in the adjunction morphism $u_!u^*F \longrightarrow F$ in $I^{\wedge}$ are functors $I^{\wedge} \longrightarrow I^{\wedge}$ which commute with arbitrary $\varinjlim$. Thus to check that it is \texttt{[an]} isomorphism, it suffices to check it for $F$ in $I$, since $I$ generates $I^{\wedge}$ by $\varinjlim$ . On the other hand, an argument of this type does not seem to prove that to check that we have an \textit{isomorphism} $F \longrightarrow u_*u^*F$, it suffices to do so for representable $F$, i.e. to check b$''$ for representable $F$.

It would remain to see whether c$'$ $\Longrightarrow$ d$''$ \texttt{[c$'$?]}.

\msp{page 55}Let us make explicit the adjunction morphism
\[
i \longrightarrow u_*u^*i
\]
(for $i \in \Ob I$) by its values on $j$
\[
\underbrace{i(j)}_{\Hom(j,i)} \longrightarrow \Gamma(I'/j, u^*i). \tag*{$(\ast)$}
\]
So let
\[
u_j : I'/j \longrightarrow I,
\]
% foot: Version 8 juillet 2001
\printedpage{35}
% head: A. Grothendieck : Dérivateurs, Ch. II, Édition des Universités de Montpellier II et Paris VII
we have
\[
\begin{aligned}
\Gamma(I'/j, \underbrace{u^*(i)|I'/j}_{u_j^*(i)}) &\simeq \Hom(e_{I'/j}, u^*(i)) \\ % ?? uncertain: extent of the underbrace -- over "u^*(i)|I'/j" or including the closing parenthesis
&\overset{\mathrm{adj}}{\simeq} \Hom(u_{j,!}e_{I'/j}, i),
\end{aligned}
\]
and the arrow $(\ast)$, functorial in $i$, is deduced from an arrow
\[
u_{j,!}(e_{I'/j}) \longrightarrow j. \tag*{$(\ast\ast)$}
\]
We see at once, as above, that $(\ast\ast)$ is \texttt{[an]} isomorphism if and only if for every $\alpha : k \longrightarrow j$, the category $k\backslash\underset{\alpha}{I'}/j$ is 0-connected, and if we require this for every $j$, $k$, $\alpha$, this is condition d$''$. But it is not clear that for $(\ast\ast)$ to be an isomorphism, it suffices that the Hom of the two sides with values in every $i$ in $I$ (and not every $F$ in $I^{\wedge}$) be \texttt{[an]} isomorphism. The difficulty comes from the fact that $u_{j,!}(e_{I'/j})$ is not necessarily representable, and that an arbitrary $F$ in $I^{\wedge}$ \msp{page 56}is not necessarily $\varprojlim$ of representable functors. So it is doubtful whether c$'$ \texttt{[c$'$?]} suffices to imply that $u$ is $W_0$-connected.

\textbf{Corollary 4.2.} \textit{For $u$ to be $W_0$-connected, it} suffices \textit{that it be a universal $W_0$-equivalence, i.e. (cf. 3.1) that the fibers be \textup{0}-connected, and that for every nonidentity arrow in $I$, defining $\bDelta^1 \longrightarrow I$, $I' \times_I \bDelta^1$ be \textup{0}-connected.}

4.3. This condition of universal $W_0$-equivalence is however \textit{strictly} stronger than $W_0$-connectedness, as we see from the example where $u$ is an \textit{equivalence of categories} not surjective on objects (e.g. the inclusion of a vertex in a 1-connected groupoid having more than one object). More generally, every localization functor is $W_0$-connected, whereas its fibers are not necessarily connected, and a fortiori it is not a universal $W_0$-equivalence.

However, if every endomorphism $p$ of an $i$ in $I$, such that $p^2 = p$ (projector) is the identity, and if \msp{page 57}$I' \longrightarrow I$ is ``transportable'', e.g. if in $I$ every isomorphism is an identity, then $u$ $W_0$-connected implies that the fibers of $u$ are 0-connected, as we see by expressing d$''$ for $\alpha$ an identity arrow. The conditions envisaged on $I$ are satisfied if $I$ is ordered. But even if $I'$ and $I$ are both ordered, $u$ can be $W_0$-connected, without being a universal $W_0$-equivalence. Here is an example with $I = \bDelta^2$:
\[
\begin{array}{r@{\;}l@{\quad}c}
I' = & \left\{\begin{tikzcd}[ampersand replacement=\&]
 \& 1' \arrow[r] \arrow[d] \& 2 \\
0 \arrow[r] \& 1 \&
\end{tikzcd}\right\} & \\[1ex]
 & & \Big\downarrow \\[1ex]
I = & \left\{\begin{tikzcd}[ampersand replacement=\&]
0 \arrow[r] \& 1 \arrow[r] \& 2
\end{tikzcd}\right\}. &
\end{array}
\]
% foot: Version 8 juillet 2001

\printedpage{36}
% head: A. Grothendieck : Dérivateurs, Ch. II, Édition des Universités de Montpellier II et Paris VII
The bilocalized categories of $I$ are the three vertices $\{0\}$, $\{1\}$, $\{2\}$, and the three segments $[0,1]$, $[1,2]$ and $[0,2] = I$, the inverse images are respectively
\[
\bullet\ ,\quad
\begin{tikzcd}[row sep=normal] {} \arrow[d] \\ {} \end{tikzcd}\ ,\quad
\bullet\ ,\quad
\begin{tikzcd}[row sep=normal] & {} \arrow[d] \\ {} \arrow[r] & {} \end{tikzcd}\ ,\quad
\begin{tikzcd}[row sep=normal] {} \arrow[r] \arrow[d] & {} \\ {} & \end{tikzcd}\ ,\quad
I'\quad ,
\]
they are all 0-connected, i.e. $W_0$-aspheric, and even contractible hence $W_\omega$-aspheric (so $u$ is even $W_\omega$-connected, i.e. $W_\omega$-biaspheric). But it is not a universal $W_0$-equivalence, since the inverse image of $I'$ under $\bDelta^1 \xrightarrow{\partial} \bDelta^2$, $\partial(0) = 0$, $\partial(1) = 2$, \msp{page 58}is the disconnected category $\varepsilon = \{\bullet\ \ \bullet\}$.

I had hoped for a moment that the $W_0$-connected functors $I' \xrightarrow{u} I$ would be those which factor as $I' \xrightarrow{u_0} I'' \xrightarrow{v} I$, with $u_0$ an equivalence, and $v$ a universal $W_0$-equivalence. But the preceding example defeats this, with $u$ even $W_\omega$-connected. Indeed, $I''$ would be derived from a preordered set, and $v$ would factor through the corresponding ordered set, namely $I'$, hence would be of the form
\[
I'' \longrightarrow I' \xrightarrow{u} I,
\]
where $I'' \longrightarrow I'$ is a passage to the quotient, and the inverse image of $\bDelta^1 \xrightarrow{\partial} I$ would again be disconnected (a $\pi_0$ of cardinality 2). So it is really hopeless.

However

\textbf{Corollary 4.3.} \textit{Suppose that $u$ is $W_0$-smooth, or $W_0$-proper, or a weak $W_0$-fibration. Then $u$ is $W_0$-connected if and only if it is a universal $W_0$-equivalence, i.e. if and only if a) its fibers are \textup{0}-connected and b) $I' \times_I \bDelta^1$} \texttt{[is]} \textit{\textup{0}-connected for every $\bDelta^1 \longrightarrow I$ (corresponding to a nonidentity arrow).}

\msp{page 59}In all three cases we know that $u$ is a universal $W_0$-equivalence if and only if its fibers are $W_0$-connected. Now $W_0$-smoothness implies that $\forall i \in \Ob I$, the inclusion
\[
I'_i \longrightarrow i\backslash I'
\]
is a $W_0$-equivalence, on the other hand $i\backslash I' \longrightarrow i\backslash I$ being $W_0$-aspheric, hence a $W_0$-equivalence by hypothesis, $i\backslash I'$ is 0-connected, hence $I'_i$ also. The case $u$ $W_0$-proper is dual. In the case \texttt{[where]} $u$ \texttt{[is]} a weak $W_0$-fibration, we again know that $I'_i \longrightarrow i\backslash I'$ is a $W_0$-equivalence, and we conclude as in the smooth case (or, as one chooses, via $I'_i \longrightarrow I'/i$).

But it occurs to me that the $W_0$-connected functors might be exactly the localization functors? It amounts to the same thing to say that a \textit{conservative} functor $u : I' \longrightarrow I$ (i.e. $u(\alpha)$ isomorphism $\Longrightarrow$ $\alpha$ isomorphism) which is $W_0$-connected, is an equivalence. But this is false, as we see by taking $I = e$, $I'$ a 0-connected and not 1-connected groupoid: $I' \longrightarrow e$ is $W_0$-connected (and even a universal $W_0$-equivalence), but it is not an equivalence (and not a localization functor).
% foot: Version 8 juillet 2001

\printedpage{37}
% head: A. Grothendieck : Dérivateurs, Ch. II, Édition des Universités de Montpellier II et Paris VII
\msp{page 60}I would like to generalize the notion of a $W_0$-biaspheric arrow to the case of an arbitrary derivator $\bD$.

\textbf{Proposition 4.4.} \textit{Let $\bD$ be a derivator on} Diag, \textit{and $u : I' \longrightarrow I$ an arrow in} Diag. \textit{Suppose that $u_*$ and $u_!$ exist, and are computed in the standard way, for every arrow $u$ in} Diag. \textit{(In fact, it suffices that this be the case for $u$ and its localized arrows $u/i : I'/i \longrightarrow I/i$, or $u$ and its colocalized arrows $i\backslash u : i\backslash I' \longrightarrow i\backslash I$.) Consider the conditions:} % [sic] pour tout flèche -> pour toute flèche, gender agreement
\begin{itemize}[leftmargin=4em]
\item[\textit{a)}] $u^* : \bD(I) \longrightarrow \bD(I')$ \texttt{[is]} \textit{fully faithful.}
\item[\textit{a$'$)}] \textit{$\forall$ $F$ in $\bD(I)$, the arrow}
\[
u_!u^*F \longrightarrow F
\]
\textit{is an isomorphism.}
\item[\textit{a$''$)}] \textit{$\forall$ $F$ in $\bD(I)$, the arrow}
\[
u_*u^*F \longleftarrow F
\]
\textit{is an isomorphism.}
\item[\textit{b)}] \textit{$\forall$ $i$ in $I$, the arrow $i\backslash I' \longrightarrow i\backslash I$ is $\bD$-aspheric.}
\item[\textit{b$'$)}] \textit{$\forall$ $i$ in $I$, the arrow $I'/i \longrightarrow I/i$ is $\bD$-coaspheric.}
\item[\textit{b$''$)}] \textit{$\forall$ $\alpha : i_0 \longrightarrow i_1$ in $I$, the category $i_0\backslash\underset{\alpha}{I'}/i_1$ is $\bD$-aspheric.}
\end{itemize}
\textit{We then have the implications}
\[
(b \Longleftrightarrow b' \Longleftrightarrow b'') \Longrightarrow (a \Longleftrightarrow a' \Longleftrightarrow a'').
\]
\textit{We have $a \Longrightarrow b$ (i.e. all the legitimate conditions envisaged are equivalent)} \msp{page 61}\textit{in each of the following cases.}
\begin{itemize}[leftmargin=4em]
\item[$1^o$] \textit{$I$ is preordered.}
\item[$2^o$] \textit{$\forall$ $i$ in $I$, $\alpha_i^* : \bD(I) \longrightarrow \bD(I/i)$ is essentially surjective.}
\item[$3^o$] \textit{$\forall$ $i$ in $I$, considering $\alpha_i : I/i \longrightarrow I$, the functor $\alpha_{i,!} : \bD(I/i) \longrightarrow \bD(I)$ is conservative.}
\item[$4^o$] $\mathcal{D} = \mathrm{HOT}_{W_0} \approx \bD_{\mathrm{Ens}}$.
\end{itemize}
\textsc{Proof}. The equivalence of a, a$'$, a$''$ is clear, say $A$. Likewise the equivalence of b, b$'$, b$''$, say $B$. % [sic] l'equivalence -> l'équivalence, accent missing

Let us prove
\[
B \Longrightarrow A
\]
assuming e.g. the existence of $u_*$. We tacitly assume, along with the existence of $u_*$ : $\bD(I') \longrightarrow \bD(I)$, the method of computing the fibers of $u_*(-)$. Thus $A$, via a$''$, is equivalent to
\[
H^{\bullet}_{\bD}(I'/i, u^*F) \xleftarrow{\mathit{iso}} H^{\bullet}_{\bD}(I/i, F) \simeq F(i)
\]
% foot: Version 8 juillet 2001
\printedpage{38}
% head: A. Grothendieck : Dérivateurs, Ch. II, Édition des Universités de Montpellier II et Paris VII
for every $F$ in $\bD(I)$, every $i$ in $I$.

Let
\[
\alpha_i : I/i \longrightarrow F, \qquad F_i = \alpha_i^*(F), % [sic] I/i \longrightarrow F -> I/i \longrightarrow I, alpha_i is defined as I/i -> I in 3^o
\]
we must therefore express that for every $F_i$ coming from an $F$, we have
\[
H^{\bullet}_{\bD}(I/i, F_i) \xrightarrow{\mathrm{iso}} H^{\bullet}_{\bD}(I'/i, (u/i)^*(F_i)). \tag*{$(\ast)$}
\]
Now subject to the existence of $(u/i)_!$, to say that $(\ast)$ is \texttt{[an]} isomorphism for \textit{every} $F_i$ in $\bD(I/i)$, means that
\[
u/i : I'/i \longrightarrow I/i
\]
\msp{page 62}is $\bD$-coaspheric, whence b$'$ $\Longrightarrow$ a$''$. It is clear that the reverse implication holds if $\alpha_i^*$ is essentially surjective, this is case $2^o$. Now $2^o$ is satisfied if $I$ is preordered, whence case $1^o$. Case $4^o$ has already been treated (4.1). We must still treat case $3^o$, which moreover includes case $2^o$, hence case $1^o$ (and also case $4^o$, cf. below).

We must express the isomorphy of the arrow $(\ast)$ for $F_i$ of the form $\alpha_i^*(F)$. Let $\xi$ be \texttt{[an]} object of $\mathcal{A} = \mathcal{D}(e)$, we must express the isomorphy of $\Hom(\xi, -)$ in the two sides of $(\ast)$, i.e. of
\[
\begin{aligned}
\Hom(I/i, F_i) \quad &\longrightarrow \quad \Hom(\xi'_{I'/i}, (u/i)^*(F_i)) \\
&\overset{\text{adj.}}{\simeq} \quad \Hom((u/i)_!(\xi'_{I'/i}), F_i),
\end{aligned}
\tag*{$(\ast\ast)$}
\]
\texttt{[rather} $\xi_{I'/i}$\texttt{]} arrow which is deduced from
\[
(u/i)_!(u/i)^*(\xi_{I'/i}) \xrightarrow{\text{adj.}} \xi_{I'/i} \tag*{$(\ast\ast\ast)$}
\]
by taking the Hom into $F_i$. The condition $u/i$ coaspheric (i.e. b$'$) means that $(\ast\ast\ast)$ is an isomorphism for every $\xi$ in $\mathcal{A}$. Does it \textit{suffice} for this that its Hom into an $F_i$ be bijective, when we restrict $F_i$ in $\bD(I/i)$ to be in the image of $\bD(I)$ under $\alpha_i^*$? Now this arrow $(\ast\ast)$, for $F_i = \alpha_i^*F$, is made explicit by the adjunction formula relative \msp{page 63}to the functor $\alpha_i$, as
\[
\Hom(\alpha_{i,!}(\xi_{I/i}), F) \longrightarrow \Hom(\alpha_{i,!}(\ldots), F).
\]
To say that it is an isomorphism for every $F$, thus means \textit{exactly} that the transform of the arrow $(\ast\ast\ast)$ under $\alpha_!$ is an isomorphism. To conclude from this that $(\ast\ast\ast)$ itself is \texttt{[an]} isomorphism, the natural condition to impose is therefore that $\alpha_{i,!} : \bD(I/i) \longrightarrow \bD(I)$ be conservative. This is case $3^o$. This completes the proof of 4.4.

4.5. Let us show that condition $3^o$ on $I$ is not always satisfied, in the case of $\bD = (\mathrm{HOT}_{W_0})^{\mathrm{opp}} \approx (\bD_{\mathrm{Ens}})^{\mathrm{opp}}$, although in this case, it is nevertheless true that for every $I'$ over
% foot: Version 8 juillet 2001

\end{document}
